\documentclass[12pt]{article}
\usepackage[utf8]{inputenc}
\usepackage[T1]{fontenc}
\usepackage{hyperref}
\title{Uneme Language Resources for K-12 Curriculum Development}
\author{Compiled by Claude Sonnet 5 (claude-sonnet-5) for Andr\'as Kornai}
\date{}
\begin{document}
\maketitle
\section{Introduction}
Uneme (ISO 639-3 une; Glottocode unem1238) is an Edoid language of Edo State, Nigeria, with about 20,000 speakers cited for 2000. My search found no Uneme-specific primers, textbooks, dictionaries or orthography. The material below is background only.
\section{Primary Educational Resources}
None found for Uneme.
\section{Linguistic and Reference Materials}
Overview article: \url{https://en.wikipedia.org/wiki/Uneme_language}. Materials on the neighbouring Edo (Bini) language, not Uneme: an Edo orthography paper (\url{https://www.centreforedostudies.be/Papers/Omozuwa,%20Victor%20Edosa%3B%20Evbuomwa,%20Osaigbovo%20Obed%20-%20Edo%20Orthography_%20The%20Compelling%20Need%20for%20Harmonization%20and%20Standardization.pdf}) and Munro's English--Edo wordlist (\url{https://centreforedostudies.be/Munro/Munro%20-%20archive_org.pdf}).
\section{Religious and Community Resources}
A Joshua Project profile exists (\url{https://joshuaproject.net/index.php/people_groups/15713/NI}). I did not verify any Uneme scripture or community organisation.
\section{Digital and Audio Resources}
None found.
\section{Curriculum Development Context}
The language is not known to be taught in schools, and no Uneme teaching materials turned up. A curriculum would have to start from new documentation and orthography work. The University of Benin repository and fieldwork-based theses may be worth checking.
\end{document}